\documentclass[11pt,a4paper]{article}
\usepackage[margin=2.4cm]{geometry}
\usepackage{amsmath,amsthm,mathtools}
\usepackage{fontspec}
\usepackage{unicode-math}
\setmainfont{STIX2Text}[Path=fonts/, Extension=.otf,
  UprightFont=*-Regular, ItalicFont=*-Italic,
  BoldFont=*-Bold, BoldItalicFont=*-BoldItalic]
\setmathfont{STIX2Math}[Path=fonts/, Extension=.otf]
\usepackage[protrusion=true,expansion=false]{microtype}
\usepackage{booktabs,enumitem}
\usepackage{placeins}
\usepackage{tikz}
\usetikzlibrary{arrows.meta,positioning,calc,fit,backgrounds}
\definecolor{tierspend}{RGB}{200,60,60}
\definecolor{tierfull}{RGB}{210,140,40}
\definecolor{tierin}{RGB}{60,130,180}
\definecolor{tierrecv}{RGB}{70,150,90}
\usepackage[psdextra,hidelinks]{hyperref}

\emergencystretch=3em
\hbadness=10000
\hfuzz=2pt

\newtheorem{theorem}{Theorem}[section]
\newtheorem{lemma}[theorem]{Lemma}
\newtheorem{proposition}[theorem]{Proposition}
\newtheorem{corollary}[theorem]{Corollary}
\theoremstyle{definition}
\newtheorem{definition}[theorem]{Definition}
\newtheorem{construction}[theorem]{Construction}
\newtheorem{assumption}[theorem]{Assumption}
\newtheorem{remark}[theorem]{Remark}
\newtheorem{example}[theorem]{Example}

% kind-coded theorem blocks, palette shared with the figures and the website:
% definitions/constructions blue (tierin), assumptions amber (tierfull),
% proved statements a neutral bar; remarks and examples run unframed
\usepackage{mdframed}
\providecolor{tierin}{RGB}{60,130,180}
\providecolor{tierfull}{RGB}{210,140,40}
\mdfdefinestyle{kindbar}{topline=false,bottomline=false,rightline=false,
  linewidth=1.5pt,innerleftmargin=9pt,innerrightmargin=4pt,
  innertopmargin=5pt,innerbottommargin=6pt,leftmargin=0pt,rightmargin=0pt,
  skipabove=10pt,skipbelow=10pt}
\surroundwithmdframed[style=kindbar,linecolor=tierin!70,backgroundcolor=tierin!4]{definition}
\surroundwithmdframed[style=kindbar,linecolor=tierin!70,backgroundcolor=tierin!4]{construction}
\surroundwithmdframed[style=kindbar,linecolor=tierfull!80,backgroundcolor=tierfull!5]{assumption}
\surroundwithmdframed[style=kindbar,linecolor=black!30]{theorem}
\surroundwithmdframed[style=kindbar,linecolor=black!30]{lemma}
\surroundwithmdframed[style=kindbar,linecolor=black!30]{proposition}
\surroundwithmdframed[style=kindbar,linecolor=black!30]{corollary}

\usepackage{fancyhdr}
\pagestyle{fancy}
\fancyhf{}
\fancyhead[L]{\footnotesize\scshape The Zcash Arboretum}
\fancyhead[R]{\footnotesize\itshape\nouppercase{\leftmark}}
\fancyfoot[C]{\footnotesize\thepage}
\renewcommand{\headrulewidth}{0.3pt}
\renewcommand{\sectionmark}[1]{\markboth{\thesection\ \ #1}{}}

\newcommand{\F}{\mathbb{F}}
\newcommand{\Z}{\mathbb{Z}}
\newcommand{\N}{\mathbb{N}}
\newcommand{\G}{\mathbb{G}}
\newcommand{\E}{\mathbb{E}}
\newcommand{\PP}{\mathbb{P}}
\renewcommand{\Pr}{\mathrm{Pr}}
\newcommand{\Fp}{\F_p}
\newcommand{\Fq}{\F_q}
\newcommand{\ip}[2]{\langle #1,\, #2 \rangle}
\newcommand{\Comm}{\mathsf{Commit}}
\newcommand{\Open}{\mathsf{Open}}
\newcommand{\Setup}{\mathsf{Setup}}
\newcommand{\Verify}{\mathsf{Verify}}
\newcommand{\negl}{\mathsf{negl}}
\newcommand{\poly}{\mathsf{poly}}
\newcommand{\PPT}{\mathsf{PPT}}
\newcommand{\Adv}{\mathcal{A}}
\newcommand{\Prov}{\mathcal{P}}
\newcommand{\Ver}{\mathcal{V}}
\newcommand{\Sim}{\mathcal{S}}
\newcommand{\Ext}{\mathcal{E}}
\newcommand{\Rel}{\mathcal{R}}
\newcommand{\Lang}{\mathcal{L}}
\newcommand{\nfsym}{\mathsf{nf}}
\newcommand{\cmsym}{\mathsf{cm}}
\newcommand{\cmx}{\mathsf{cmx}}
\newcommand{\cvsym}{\mathsf{cv}}
\newcommand{\rk}{\mathsf{rk}}
\newcommand{\Pallas}{\ensuremath{\mathsf{Pallas}}}
\newcommand{\Vesta}{\ensuremath{\mathsf{Vesta}}}
\newcommand{\Ep}{\ensuremath{\mathcal{E}}}
\newcommand{\NoteCommit}{\ensuremath{\mathsf{NoteCommit}}}
\newcommand{\Sinsemilla}{\ensuremath{\mathsf{Sinsemilla}}}
\newcommand{\ShortCommit}{\ensuremath{\mathsf{ShortCommit}}}
\newcommand{\CommitIvk}{\ensuremath{\mathsf{Commit}^{\mathsf{ivk}}}}
\newcommand{\Extract}{\ensuremath{\mathsf{Extract}}}
\newcommand{\Acc}{\ensuremath{\mathsf{Acc}}}
\newcommand{\GroupHash}{\ensuremath{\mathsf{GroupHash}}}
\newcommand{\ask}{\ensuremath{\mathsf{ask}}}
\newcommand{\aksym}{\ensuremath{\mathsf{ak}}}
\newcommand{\nksym}{\ensuremath{\mathsf{nk}}}
\newcommand{\ivk}{\ensuremath{\mathsf{ivk}}}
\newcommand{\fvk}{\ensuremath{\mathsf{fvk}}}
\newcommand{\ovk}{\ensuremath{\mathsf{ovk}}}
\newcommand{\dk}{\ensuremath{\mathsf{dk}}}
\newcommand{\pkd}{\ensuremath{\mathsf{pk_d}}}
\newcommand{\gd}{\ensuremath{\mathsf{g_d}}}
\newcommand{\rcv}{\ensuremath{\mathsf{rcv}}}
\newcommand{\rcm}{\ensuremath{\mathsf{rcm}}}
\newcommand{\rivk}{\ensuremath{\mathsf{r_{ivk}}}}
\newcommand{\rsk}{\ensuremath{\mathsf{rsk}}}

\renewenvironment{abstract}{%
  \small\begin{center}{\bfseries\abstractname}\end{center}\medskip\noindent\ignorespaces
}{\par\medskip}

\title{\textbf{\Huge Sync Guide}\\[6pt]\large Discovering notes and maintaining spendable state}
\author{m@rek.onl}
\date{}
\begin{document}
\maketitle
\thispagestyle{empty}
\begin{abstract}
This volume of \emph{The Zcash Arboretum} constructs the state machine that
turns a stream of chain data into a usable shielded wallet. It assumes the
\emph{Crypto Guide}'s authenticated trees, the \emph{Consensus Guide}'s
branch-local state and reorganisation rules, the \emph{Ironwood Guide}'s
note decryption and membership equations, and the \emph{Wallet Guide}'s
account and transaction lifecycle. It explains what data to fetch, how
to recognise received money, how to maintain the evidence needed to spend
it, and how to undo observations when the selected history changes.
It then identifies the remaining server-trust and metadata boundaries.
It does not equate successful scanning with
independent consensus validation.
\end{abstract}
\clearpage
\tableofcontents
\clearpage

\section{What a wallet must reconstruct}\label{sec:sync-scope}

The blockchain publishes commitments, nullifiers and encrypted note data.
Their concrete constructions are in the \emph{Ironwood Guide}, ``Notes
and their commitments'', ``Nullifiers and authenticated existence'', and
``Note delivery and recovery''. In particular, the private path of sibling
hashes used to prove membership is the note's \emph{membership witness}.
A wallet needs a more useful private view: which outputs pay its keys,
what values they carry, whether they remain unspent, and which witness
will allow each to be spent. These facts arise from different operations.
Decryption discovers a note; commitment reconstruction checks its
content; ordering determines its position; nullifier matching observes
a spend; and chain validation or an explicit trust assumption establishes
the branch in which those observations have meaning.

\subsection{A complete wallet record}

For each detected note, retain enough information to reconstruct
\[
 (\mathsf{pool},\mathsf{account},\mathsf{scope},
  \mathsf{note},\mathsf{position},\mathsf{creation},
  \mathsf{nullifier},\mathsf{spend\ observation}).
\]
The note includes its recipient, value, creation identifier and seed, or
equivalent fully reconstructed components. Its creation record identifies
block hash, height, transaction identity and Action/output index.
A spend observation is branch-linked and may be absent. A note
nullifier may also be unavailable to a wallet given only incoming
viewing authority.

Per pool, retain the current commitment-tree size and enough tree data
to produce required paths at chosen checkpoints. A \emph{checkpoint} is
a retained tree state associated with a particular block boundary.
Its root and size must be bound to a branch identity, not just a height.

These records serve the \emph{Wallet Guide}, ``Fees, inputs, and change''.
A detected note without a usable witness can contribute to a displayed
balance while not yet being selectable for a transaction.

\subsection{Inputs and authority}\label{sec:sync-authority}

A sequential scanner needs:
\begin{enumerate}[leftmargin=*]
\item account viewing keys and the external/internal scopes to search;
\item a starting branch point and its pool tree states;
\item complete relevant commitment and nullifier sequences after that
  point, in their protocol order;
\item output data sufficient for note discovery;
\item a way to decide whether the supplied branch and its data are
  authentic, or an explicit trust decision about their source.
\end{enumerate}
The last item cannot be supplied by the first four. Matching an encrypted
output to one's key does not establish that the output was mined or
that the server showed every spend.

The normative and implementation reference baseline is
\texttt{zips} commit \texttt{69610984}, especially the note-decryption
rules and ZIP~307, and \texttt{librustzcash} commit
\texttt{e30517e4}, especially
\texttt{zcash\_client\_backend/src/scanning.rs},
\texttt{src/data\_api/chain.rs} within that crate, and
\texttt{proto/compact\_formats.proto}. The compact schema's inclusion of
Ironwood fields is capability evidence. It is not proof that their pool
rules have activated on a particular network.

\section{Compact discovery}\label{sec:sync-compact}

A full transaction contains proofs and signatures that a basic
key-discovery scan does not need. A \emph{compact block} is a service
representation retaining selected discovery and tree-maintenance data.
It is not another consensus block format.

\subsection{The information retained}\label{sec:sync-compact-wire}

A compact Orchard-protocol Action carries:
\begin{center}
\begin{tabular}{@{}lrp{0.50\linewidth}@{}}
\toprule
Field & Bytes & Purpose\\
\midrule
Nullifier & 32 & Spend detection; creation identifier of the paired output.\\
Extracted commitment & 32 & Output reconstruction and tree insertion.\\
Ephemeral key & 32 & Incoming key agreement.\\
Ciphertext prefix & 52 & Lead byte, diversifier, value and note seed.\\
\bottomrule
\end{tabular}
\end{center}
These fields total $148$ payload bytes before service encoding overhead.
The omitted incoming ciphertext contains $512$ encrypted memo bytes and
the $16$-byte authentication tag. The $80$-byte outgoing ciphertext,
net value commitment, randomised key, proof and signatures are also not
part of this compact Action.

A compact transaction supplies its original transaction index and
identifier, plus the relevant pool lists. A compact block supplies height,
hash, parent hash and pool tree-size metadata where supported.
Transaction indices can have gaps when the service omits transactions
with no requested relevant components. The retained lists still have to
preserve the exact order of commitments and nullifiers.

The specified Ironwood extension uses a distinct Action list and tree
size. Its Action payload has the same shape. Do not merge its commitments
with the original Orchard list: the two pools have separate position
spaces and roots even though their receivers and primitive hashes are
shared.

The service uses \emph{Protocol Buffers}, a message encoding with numbered
fields and prescribed default values. It can represent an absent field by its default zero or
empty value. A scanner must interpret presence and metadata semantics
according to the service contract; an omitted tree size is not
automatically an authenticated assertion that the pool is empty.
The compact block's optional header is not guaranteed to be populated.
The baseline schema explicitly describes normal compact headers as
empty.

\subsection{Decrypting without an authentication tag}
\label{sec:sync-compact-decrypt}

First parse the Action's field and point encodings canonically and enforce
the required ephemeral-point conditions. Only then, for each incoming
scalar $\ivk$, compute the shared point
$[\ivk]\mathsf{epk}$ and the KDF from the \emph{Ironwood Guide},
``Note delivery and recovery''. The first $52$ ciphertext bytes use the
same ChaCha20 keystream prefix as the full encrypted note. Invert that
prefix to obtain a candidate
\[
 (\ell,d,v,\mathsf{rseed}).
\]
There is no Poly1305 tag in the compact data, so this is not a successful
authenticated-encryption-with-associated-data (AEAD) opening.
Correctness is tested by reconstructing the note:
\begin{enumerate}[leftmargin=*]
\item Require canonical outer field and point encodings and the
  appropriate ephemeral-point conditions.
\item Check that lead byte $\ell$ is allowed for the pool's note format.
\item Set $\rho$ to this Action's nullifier. Derive $\gd$ from $d$ and
  $\pkd=[\ivk]\gd$.
\item Derive $\mathsf{esk}$ and $\psi$, then the correct version's
  $\rcm$. Require the derived ephemeral point to match the supplied one.
\item Recompute the commitment and compare its extraction with the
  Action's $\cmx$.
\end{enumerate}
Only a candidate satisfying those checks becomes a detected note.
The security argument is the note-reconstruction contract, not the
missing AEAD tag. When the wallet later retrieves the full transaction,
it performs full ciphertext authentication and checks that the
retrieved output belongs to the same committed transaction context.

A malformed compact record is not a well-formed record that merely
decrypts to somebody else's key. A production scanner must distinguish
structural failure from a harmless no-match result. It must not silently
omit a malformed commitment and shift every subsequent note position.
If the service does not supply trustworthy complete data, the safe
response is to reject or repair the range, not mark it completely scanned.

\subsection{Incoming discovery and spend discovery}

For $k$ incoming scopes and $m$ output records, straightforward discovery
performs $km$ candidate key agreements and reconstruction attempts.
Batching and prepared-key operations can reduce cost without changing
which key/output pairs must be considered. Deriving many diversified
addresses under one scope does not multiply $k$.

For an account with a full viewing key, a detected note's nullifier is
computed by the \emph{Ironwood Guide}, ``Nullifiers and authenticated
existence''. Match that value against all observed nullifiers in the same
pool. A match means the note was consumed on the supplied branch.
Its Action pairing does not tell a recipient which input financed its
output: the incoming nullifier also serves as the output's creation
identifier, but the old commitment remains private.

An incoming-only key lacks the nullifier derivation capability. It can
discover received value without independently reconstructing ordinary
spentness for those notes. Exporting an incoming key and calling the
result a complete unspent balance would promise more than that
capability permits.

\section{Commitment positions and witnesses}\label{sec:sync-tree}

The proof needs a path to an eligible root. The encrypted note does not
contain that path or its position. They arise from the ordered public
commitment history.

\subsection{Positions are arithmetic}\label{sec:sync-tree-linear}

Suppose a pool contains $n$ commitments before a block, and the block
appends $c$ commitments. Their positions are
\[
 n,n+1,\ldots,n+c-1,\qquad n_{\rm after}=n+c.
\]
Every Action contributes its output commitment, including a zero-valued
dummy output. Counting only notes detected by this wallet is incorrect.

For a transaction with several Actions, add the count from all previous
transactions' same-pool Actions to $n$, then add the Action's within-list
index. If reliable end-of-block metadata gives $n_{\rm after}$, one can
derive $n=n_{\rm after}-c$, checking underflow and continuity with known
states. This arithmetic locates leaves; it does not authenticate the
metadata or prove that $c$ includes every required commitment.

For the shared payment, suppose the pool had $11$ leaves before the
transaction. Its payment and change occupy positions $11$ and $12$ in
the chosen Action order, even if Bob discovers only the first and Alice
only the second. A preceding transaction's commitments would shift both
positions by its count. A reorganisation can change these positions.

\subsection{Empty roots and the incremental frontier}
\label{sec:sync-tree-frontier}

Use the concrete height-indexed hash $H_j$ from the \emph{Ironwood Guide},
``Nullifiers and authenticated existence''. Let $e_0=2$ be the empty-leaf
sentinel and define
\[
 e_{j+1}=H_j(e_j,e_j)\quad(0\leq j<32).
\]
These are the roots of completely empty subtrees of increasing height.

A \emph{frontier} stores the roots of completed left subtrees needed to
continue appending. For $n<2^{32}$ leaves, its occupied heights are the
set bits of $n$. At height $j$, the retained value $F_j$ covers the
corresponding completed subtree in the ordered prefix.
This is binary carrying with subtree hashes instead of digits.

To append a leaf $L$, first require $n<2^{32}$. Set $x=L,j=0$ and
$t=n$. While $t$ is odd, combine
$x\leftarrow H_j(F_j,x)$, remove $F_j$, set
$t\leftarrow\lfloor t/2\rfloor$ and increment $j$.
When $t$ is even, store $F_j=x$ and increment the original count.
At the exact full capacity, retain the resulting height-$32$ root as
a terminal full-tree state; there is no next append.

\begin{example}[A binary carry]
For $n=11=8+2+1$, frontier heights are $3,1,0$.
Appending the twelfth leaf first joins the old height-zero subtree,
then joins the old height-one subtree, producing a new height-two
subtree. The frontier heights become $3,2$, matching $12=8+4$.
No earlier leaf is rehashed.
\end{example}

For a nonfull tree, recover the padded depth-$32$ root without storing
empty leaves. Start $x=e_0$. At each height $j=0,\ldots,31$, set
\[
 x\leftarrow
 \begin{cases}
 H_j(F_j,x),&\text{if bit }j\text{ of }n\text{ is }1,\\
 H_j(x,e_j),&\text{otherwise}.
 \end{cases}
\]
The current $x$ is the partial rightmost subtree of height $j$.
A set bit supplies its completed left sibling; an unset bit supplies an
empty right sibling. Induction on height proves that after the final
step $x$ is the root of the prefix padded to capacity.

The full frontier needs at most one field element per height and the
count. It is sufficient to append and compute roots.
It is \emph{not} sufficient to reconstruct an arbitrary old leaf's
membership path: hashes deliberately discard the child information.
A wallet retaining owned notes needs additional marked subtrees or
witness data.

\subsection{Retaining paths for marked leaves}\label{sec:sync-tree-witness}

Mark a leaf when its note is detected. A simple complete implementation
retains all tree nodes; it can answer a path query by retrieving the
sibling subtree root at every height. A smaller wallet retains only the
nodes needed for marked leaves and checkpoints, collapsing unrelated
completed subtrees to their roots.

For a marked position $i$ and a height $j$, its sibling interval begins
at
\[
 a_j=\bigl((\lfloor i/2^j\rfloor)\mathbin{\mathrm{xor}}1\bigr)2^j
\]
and covers $2^j$ leaves. The sibling value is the root of that interval,
padded with the appropriate empty roots where the interval extends beyond
the checkpoint's populated prefix. Return these $32$ values with $i$;
the verifier recomputes the root using the bit-directed path equation.

An append at position $p>i$ can change only the sibling subtree whose
interval contains $p$. Its height is the highest set-bit index of
$i\mathbin{\mathrm{xor}}p$. Update that subtree's partial root and retain
the other siblings. A root-only collapsed subtree entirely outside all
future queried paths needs no retained interior.

This retention policy has a limit. If a wallet later imports a new key
and discovers an old note inside a subtree whose interior was discarded,
that subtree root cannot reveal the missing path. The wallet must
retrieve the necessary earlier commitments or a correctly authenticated
witness. Small storage is achieved by choosing what future queries to
support, not by inverting hashes.

\subsection{Shards and checkpoint roots}\label{sec:sync-tree-roots}

Partition the tree into subtrees of fixed height $s$, called
\emph{shards}. Shard $j$ covers
$[j2^s,(j+1)2^s)$. A shard containing an owned note retains its required
interior. A complete shard containing no marked note can be represented
by its root. The tree above shard roots retains their ordering and
connects them to the full root.

To construct a witness, concatenate the path inside the note's shard
with the path of that shard root through the upper tree. The lengths
are $s$ and $32-s$, totalling $32$. For example, $s=16$ divides the
path into two sixteen-hash parts; it is a storage choice, not a new
consensus tree depth.

A service may provide completed subtree roots to accelerate bootstrap.
A correct root for an earlier subtree can replace replay of its interior
when that interior is not needed for a marked leaf. The wallet still
needs trustworthy placement, completion boundary and branch association.
An unverified root supplied by the same untrusted server is a shortcut
within that trust model, not an independent security proof.

A checkpoint must also retain enough of the partially filled current
shard to reproduce that block's historical root. Keeping only today's
root for a mutable partial shard cannot answer every older anchor query.
Pruning therefore follows the chosen checkpoint and spend policy:
do not remove the only information needed by a still-pending transaction
or a supported rollback boundary.

\section{Atomic scanning and rollback}\label{sec:sync-scanning}

A scan updates note records and tree state together. An \emph{atomic}
update makes all of its writes visible together or none of them.
Their consistency
is more important than whether the implementation processes one block
or many blocks per database transaction.

\subsection{The sequential state machine}\label{sec:sync-scan}

Let the current scan checkpoint be
$(h,b,\{n_P,F_P\}_P)$, with block hash $b$.
For a next compact block $B$, perform:
\begin{enumerate}[leftmargin=*]
\item Require expected height and parent hash, when these are supplied
  by the protocol, and validate all structural bounds and encodings.
  Compare against independently authenticated block information where
  available. Missing required context is not a completed scan.
\item Work in temporary state. For each pool, determine the starting
  tree size, checking consistency against prior state and the complete
  ordered commitment count.
\item Traverse compact transactions in original block order. Match
  nullifiers against tracked notes and collect spend observations.
  For every output, run candidate decryption for each relevant scope.
\item Append every commitment in pool order. A detected note receives
  exactly this append's position; mark that position and compute its
  nullifier when the supplied key capability permits.
\item Reconcile the block's collected nullifiers with detected notes,
  verify tree-size metadata, and derive the final roots and checkpoint.
  If an authenticated final root is available, compare it.
\item Atomically persist detected notes, spend observations, tree
  changes, transaction metadata and the new branch checkpoint.
  Mark the range scanned only after those writes succeed.
\end{enumerate}
A rejected block or interrupted transaction leaves the previous checkpoint
intact. Replaying an already committed range must be idempotent:
unique block/transaction/output identities prevent duplicated notes and
duplicated append effects.

The wallet's pool root after a block can be authenticated indirectly by
the later history commitment described in the \emph{Consensus Guide},
``Authenticated block history''. That path must include the correct
epoch, node metadata and header commitment. A root merely written in a
compact message is not equivalent evidence.

\subsection{Parallel and out-of-order work}\label{sec:sync-queue}

Trial decryptions for distinct key/output pairs can run independently.
Downloaded ranges can be cached and processed out of order if their tree
positions and boundary states are already known reliably.
The final database must nevertheless describe one coherent branch.

For a disconnected range $[a,b]$, do not guess its starting tree size
from the last locally scanned height. Obtain a boundary state or derive
the size from consistent authenticated metadata and exact counts.
The scanner may know that it found a note in that range before it can
produce the note's complete witness. Record this partial knowledge
honestly rather than mark the whole account fully synchronised.

A scan scheduler can prioritise recent blocks to discover recent receipts,
then fill older gaps. Its priority is local policy. The highest scanned
height is not a sufficient progress measure: an unscanned gap can hide a
note or spend. Track completed intervals and their branch associations,
and expose whether the contiguous required range is complete.

The source implementation's separation of scan ranges, note discoveries
and retained tree data reflects these different kinds of knowledge.
The explanatory algorithm does not require a particular RPC inventory
or database schema to preserve them.

\subsection{Detecting and applying a reorganisation}
\label{sec:sync-reorg}

A mismatching parent hash or a disagreement with an authenticated tip
signals that the stored branch may no longer match the selected branch.
Locate a common ancestor by comparing retained block identities with
the replacement branch. A sequence of server replies indexed only by
height must not be mistaken for proof that no reorganisation happened.

Roll back chain-derived records beyond the common ancestor:
\begin{itemize}[leftmargin=*]
\item remove or unmine notes created only on the removed suffix;
\item remove spend observations that occurred only there;
\item restore pool frontiers, marked-path state, checkpoint sets and sizes;
\item remove stale scanned-range assertions and mined transaction
  confirmations for that suffix.
\end{itemize}
Keep the user's payment intentions and retained signed transactions
separate. Re-evaluate them under the \emph{Wallet Guide}, ``Tracking
payment state and recovery''. A removed spend may become pending again,
but can also conflict or expire on the replacement branch.

\begin{example}[Positions after rollback]
Suppose a common ancestor has $11$ commitments. One branch appends the
two outputs of the shared payment at positions $11$ and $12$.
The replacement branch first appends three other commitments.
If the same payment is then included, its outputs occupy positions
$14$ and $15$. The note data and commitment can be unchanged while
the position and path change. Reusing the old path would be incorrect.
\end{example}

If the common ancestor predates retained rollback information, restart
from an earlier trusted checkpoint and rescan. A retention window is a
resource decision. It does not justify continuing with a mixed or
unrecoverable state, and it is not a proof of consensus finality.

\subsection{Restoring from a birthday}\label{sec:sync-birthday}

A birthday says where a required scan begins. If a restored account was
created before height $b$, initialise its starting tree state at the
chosen preceding boundary and scan all required data from there.
The birthday does not reconstruct that initial root or frontier; the
wallet must obtain it by replay, an authenticated checkpoint, or a
declared trusted source.

A birthday later than a received note can miss the note entirely.
A birthday earlier than necessary costs time and bandwidth.
Neither fact reveals whether a server omitted a relevant transaction
inside the scanned range. Account enumeration and restored incoming
scopes must also match the original wallet; correct block scanning under
the wrong keys cannot discover the missing account.

\section{Server trust and observable metadata}\label{sec:sync-privacy}

Running decryption locally prevents a server from learning viewing keys
merely to provide compact data. It does not make the entire interaction
private or independently validated.

\subsection{Four independent checks}\label{sec:sync-privacy-model}

Separate these questions:
\begin{enumerate}[leftmargin=*]
\item \emph{Content consistency:} does a recovered note match a supplied
  commitment and ephemeral key?
\item \emph{Inclusion:} do the supplied transaction and state data belong
  to an authenticated block or tree root?
\item \emph{Chain validity and selection:} does that block belong to the
  selected valid most-work branch under the client's model?
\item \emph{Completeness and availability:} did the client receive all
  data necessary to discover its notes and spends?
\end{enumerate}
A commitment check answers the first. A Merkle path against a trusted root
can answer a specified inclusion question. Header work alone does not
validate hidden transaction proofs or account for unavailable bodies.
Even correct individual inclusion proofs do not prove that the server
has revealed every relevant item.

A conventional lightweight scan can trust a service for full validation
and completeness while keeping keys local. State that arrangement
plainly. A full-node-backed wallet can delegate those same claims to a
locally controlled validator. An independently authenticating client
needs additional verified data and a chain-selection argument.
These are different trust configurations, not different levels of
confidence in the same unauthenticated reply.

\subsection{What requests reveal}\label{sec:sync-privacy-leaks}

A server can observe source network identifiers, connection times,
requested heights, range sizes, pauses and retries. A narrow birthday
range can reveal approximate wallet creation or restoration activity.
A follow-up request for a full transaction immediately after compact
scanning can indicate that the wallet found it relevant.
A broadcast request links the submitted transaction to the requesting
connection unless another communication layer prevents that observation.

The chain itself also exposes transaction sizes, Action counts, pools,
public value balances and timing. Coin selection or unusual change
shapes can create additional correlations. Hiding note values from the
server's compact scan does not erase those public facts.

Transparent address-specific queries reveal the queried addresses to
the service. A design that avoids sending shielded viewing keys but
uploads transparent receiver lists has a different privacy boundary for
the transparent part of the account.

\subsection{Mitigations and their limits}

Larger uniform range requests can reduce height-resolution leakage.
Batching full-transaction retrieval can obscure which compact match
triggered a request. Separating scan and broadcast services can reduce
one service's direct view of both activities. Using multiple sources can
reveal disagreements or reduce reliance on one operator.
These techniques do not automatically provide anonymity: operators may
collude, network observers may correlate timing, and all sources may
repeat the same omitted data.

Fetching every full transaction avoids relevance-triggered follow-up
requests but increases bandwidth. Running one's own validating node
changes the server-trust boundary but not public-chain leakage.
The correct choice depends on the wallet's explicit bandwidth, storage
and adversary model. No mechanism in this chapter proves a universal
network privacy theorem.

\subsection{The synchronisation contract}

A wallet is ready to construct a particular spend when the required
account and scope have been scanned over the necessary coherent branch
range, the input note has a verified reconstruction and unspent status
under the chosen trust model, and its witness is available for the
selected anchor. This is more precise than saying that the latest
displayed height equals the server's tip.

The core invariant is a correspondence:
\[
 \text{recorded note}\ \longleftrightarrow\
 \text{pool and branch position}\ \longleftrightarrow\
 \text{path to the selected anchor}.
\]
Spend observations and transaction lifecycle records must refer to that
same branch. If any link is missing, retain a partial state or rescan;
do not silently convert uncertainty into spendable balance.

\paragraph{Evidence and executable checks.}
The compact schema is
\path{zcash_client_backend/proto/compact_formats.proto}.
The note scan and persistence interfaces in that crate are
\path{src/scanning.rs} and \path{src/data_api/chain.rs}.
The database implementation is \path{zcash_client_sqlite/src/chain.rs},
at the pinned library commit.
The generic retained-tree implementation is the
\texttt{shardtree} crate, used through those interfaces.
The companion \texttt{check\_deployed.py} checks append occupancy,
frontier roots against complete trees, every retained path and its changed
sibling through all $32$ appends of a depth-five toy tree, position
arithmetic and rollback using a clearly labelled toy hash.
It does not replace the production Sinsemilla vectors, test a deployed
server for omission, or establish an independent light-client security
theorem.
\end{document}
